\documentclass[11pt,a4paper]{article}
\usepackage[utf8]{inputenc}
\usepackage[top=0.9in, bottom=0.8in, left=1in, right=1in]{geometry}
\usepackage{hyperref}

\setlength{\parskip}{0.8em}
\setlength{\parindent}{0pt}

\begin{document}

\begin{center}
    {\LARGE \textbf{Liuxi Mei}} \\
    \vspace{2mm}
    Linköping, Sweden $\cdot$ liume102@student.liu.se $\cdot$ +46-0704898962
    \vspace{4mm}
\end{center}

\textbf{Date:} April 6, 2026 \\
\textbf{To:} Björn Forsberg \\
\textbf{Subject:} Application for PhD student in Computational Cryo-EM and Structural Biology \\

Dear Professor Forsberg,

I am writing to apply for the PhD position in Computational Cryo-EM and Structural Biology. I am finishing my master's degree right now. I will graduate and get my diploma in June this year.

I have over five years of work experience in the software industry. I also have recent research experience in machine learning. In my software jobs, I managed huge amounts of real-time data using Python. These engineering skills helped me a lot in my master's studies. For my thesis, my main job was to process very large biological datasets. I used Graph Neural Networks (GNNs) on LINCS L1000 gene expression data and OmniPath network data to study drug responses.

I do not have direct experience in cryo-EM or 3D reconstruction yet. But I have read some review papers about it, and the data pipeline is very interesting to me. I know the field now uses CNNs and Transformers to find protein particles automatically. It also uses machine learning to group 2D images and remove noise. Then, it uses deep learning to build 3D models from these 2D images. Also, we can combine AlphaFold predictions with real cryo-EM data. This makes the protein models much more accurate.

This focus on structure connects deeply to my current thesis. Biological responses happen in a clear order. First, a drug binds to a protein. Then, this triggers signaling pathways. Finally, this changes the gene expression. My current research looks at the downstream effects, like pathways and gene expression. But structural biology looks at the true upstream starting point. I really want to study this upstream level.

My career goal is to connect complex computer algorithms with real biological discoveries. My software background shows that I can handle big data. My recent thesis gives me a strong passion for data-driven biology. Your project focuses on creating new methods. This fits perfectly with my goal to build deep technical skills. Thank you for reading my letter. I hope to talk with you soon.

Sincerely,

Liuxi Mei

\vspace{0.5cm}
\noindent
\begin{minipage}{\textwidth}
\textbf{References:}

\vspace{0.3cm}
\begin{minipage}[t]{0.48\textwidth}
\textbf{Oleg Sysoev}\\
Senior Associate Professor, Linköping University\\
Department of Computer and Information Science (IDA)\\
Master's Thesis Supervisor\\
Email: \href{mailto:oleg.sysoev@liu.se}{oleg.sysoev@liu.se}
\end{minipage}%
\hfill
\begin{minipage}[t]{0.48\textwidth}
\textbf{Krzysztof Bartoszek}\\
Senior Associate Professor, Linköping University\\
Department of Computer and Information Science (IDA)\\
Bioinformatics Teacher\\
Email: \href{mailto:krzysztof.bartoszek@liu.se}{krzysztof.bartoszek@liu.se}
\end{minipage}
\end{minipage}

\end{document}